% ===========================================================================
\chapter{The cheapest way: Dijkstra's algorithm}\label{ch:dijkstra}

In 1956, Edsger Dijkstra, a young Dutch programmer, wanted a problem to
show off a new computer to people who knew nothing about computers. He
chose finding the shortest route between two cities in the Netherlands,
and, sitting at a café, worked out how to do it in twenty minutes. His
algorithm is still at the heart of every route planner.

\section{The idea}

Breadth-first search explores the frontier in the order it found it. In
Dijkstra's algorithm, every frontier node also carries the cost of the
cheapest way to it found so far, and the search always explores the
\emph{cheapest} frontier node next.

When the search explores node $n$, it looks at every edge from $n$ to a
neighbour $m$. The way to $m$ through $n$ costs the way to $n$, plus the
edge. If that is cheaper than any way to $m$ found before (or $m$ hasn't
been seen at all), it becomes $m$'s best known way: $m$ gets that cost,
and $n$ as its parent.

\begin{mathbox}[Why the cheapest frontier node is finished]
Say node $n$ is the cheapest on the frontier, with cost $c$. Could there
be a cheaper way to $n$ that the search hasn't found yet? Any such way
starts at the source, inside the explored region, and must leave the
region at some point, through some frontier node $f$. The way to $f$
already costs at least $c$, because $n$ was the cheapest on the frontier,
and the rest of the way from $f$ to $n$ can only add to that, since no edge
costs less than nothing. So no way to $n$ costs less than $c$, and $n$'s
best way is final. The search can explore it, and never look at it again.
\end{mathbox}

So the explored region grows in order of cost, like a ring of water
spreading from a stone dropped in a pond, except that the ripple runs
fast along cheap roads and slowly along expensive ones. When the target is
explored, its cost is final, and following its parents gives the cheapest
way there.

The argument needs one thing: no edge may have a negative cost. Lengths
and travel times never do.

\section{A priority queue}

``Take the cheapest frontier node'' could be done by looking through the
whole frontier every time, but there is a data structure made for it: a
\emph{priority queue}, which hands out its items smallest first. Nex has
one built in, \code{Min_Heap}. Adding an item and taking out the smallest
both take a time proportional to $\log n$ for $n$ items, instead of $n$.

What goes in the queue is a node with its cost, and its \emph{priority}:
the number the queue sorts by. For Dijkstra's algorithm the priority is
just the cost; in the next chapter it won't be. A \code{Min_Heap} of our
own class needs to be told how to compare two items: a function that gives
a negative number when the first comes first, positive when it comes
second, and 0 when they tie. \code{Real}'s \code{compare} does exactly that.

\begin{append}{search.nex}
-- A node waiting to be explored, with the cost of the best way to it
-- found so far, and its priority: the lower, the sooner it is explored.
class Waiting
create
  make(n: Integer, c, p: Real) do
    node := n
    cost := c
    priority := p
  end
feature
  once node: Integer
  once cost: Real
  once priority: Real
end

-- Lower priority first.
function by_priority(): Function(Waiting, Waiting): Integer do
  result := fn (a, b: Waiting): Integer do result := a.priority.compare(b.priority) end
end
\end{append}

A priority queue doesn't let us change the cost of an item already in it,
which we would need when a cheaper way to a waiting node turns up. Instead
we simply add the node again, with its new cost. The cheaper copy comes out
first; when the old, dearer copy comes out later, the node is already
finished, and we skip it. This is called \emph{lazy deletion}, and it is
the usual way to do it.

\section{The algorithm}

\begin{append}{search.nex}
-- Dijkstra's algorithm: the way with the lowest total cost, where `cost`
-- gives the cost of each edge (its length, or the time it takes).
--
-- Like breadth-first search, but the queue is a priority queue: the next
-- node out is always the one with the cheapest way to it found so far.
-- When a node comes out, no cheaper way to it can turn up later, because
-- any other way would have to go through a node still waiting, and every
-- one of those is already at least as expensive. So the node is finished.
function dijkstra(g: Road_Graph, source, target: Integer,
  cost: Function(Edge): Real): Search_Result
require
  source_on_roads: g.has_node(source)
  target_on_roads: g.has_node(target)
do
  let waiting: Min_Heap[Waiting] := create Min_Heap[Waiting].from_comparator(by_priority())
  waiting.insert(create Waiting.make(source, 0.0, 0.0))
  let best: Map[Integer, Real] := {source: 0.0}
  let parents: Map[Integer, Integer] := {}
  let finished: Set[Integer] := #{}
  let explored: Array[Integer] := []
  let found: Boolean := false
  from
  until
    found or waiting.is_empty
  do
    let w: Waiting := waiting.extract_min
    -- A node can be in the queue more than once, if a cheaper way to it
    -- turned up after it was first added. Only its first time out counts.
    if not finished.contains(w.node) then
      finished.add(w.node)
      explored.add(w.node)
      if w.node = target then
        found := true
      else
        across g.edges_from(w.node) as e do
          let c: Real := w.cost + cost(e)
          if not best.contains_key(e.target) or c < best.get(e.target) then
            best.put(e.target, c)
            parents.put(e.target, w.node)
            waiting.insert(create Waiting.make(e.target, c, c))
          end
        end
      end
    end
  end
  let path: Array[Integer] := []
  if found then
    path := trace_back(parents, source, target)
  end
  result := create Search_Result.make("Dijkstra", path,
    best.try_get(target, 0.0), explored, parents)
ensure
  valid_path: not result.found() or is_route(g, result.path)
  cost_adds_up: not result.found()
    or (route_cost(g, result.path, cost) - result.cost).abs < 0.001
end
\end{append}

Compare it with \code{breadth_first}: the shape is the same, a loop that
takes a node from the frontier and looks at its edges. The queue has become
a priority queue, \code{seen} has become \code{best}, which remembers a
cost for every node seen, and a node is finished when it comes \emph{out}
of the queue, not when it goes in.

The second postcondition checks the algorithm's bookkeeping against an
independent calculation: adding up the cost of each edge along the path
found must give the cost the search reports.

\section{Shortest or quickest?}

The cost function decides what ``best'' means. With \code{by_distance()},
Dijkstra's algorithm finds the shortest route; with \code{by_time()}, the
quickest.

\begin{tryit}{try_dijkstra.nex}
intern Try_Pictures

let island: Island_Map := create Island_Map.load(user_path("kestrel_island.map"))
let g: Road_Graph := create Road_Graph.build(island, Mode.Driving)
if ?island.find_place("beach") as a and ?island.find_place("lighthouse") as b then
  let source: Integer := g.nearest_node(island.point(a.node_id))
  let target: Integer := g.nearest_node(island.point(b.node_id))
  let shortest: Search_Result := dijkstra(g, source, target, by_distance())
  say("shortest  " + summary(g, shortest))
  let quickest: Search_Result := dijkstra(g, source, target, by_time())
  say("quickest  " + summary(g, quickest))
  picture_of_search(create Map_View.make(island, 1000, 700), shortest, "dijkstra.png")
end
\end{tryit}
\tryoutput{try_dijkstra}
\mappicture{dijkstra}

The shortest way goes up the middle of the island: Lake Road, then Hill
Road, then Lighthouse Lane, 3.4 km. The quickest goes round the west coast
on the Coast Road. It is 700 metres longer, but the Coast Road is a 60\,km/h
road, against 40 for the lanes, and that more than makes up for it. Real
navigators nearly always plan the quickest route, and that is what our
navigator will do too.

Breadth-first search found the quickest route as well, but only by luck:
the Coast Road happens to have both few steps and fast ones.

Look at the purple, though. Dijkstra's algorithm explored most of
the island to find a route to its northern tip, including the whole of the
town in the south, which is in exactly the wrong direction. It has no
idea where the lighthouse is; it only knows what things cost. The next
chapter gives it a sense of direction.
